% FIGURE 1 --- the capability dissociation across the three reasoning axes.
% Conceptual "money figure": a 2x3 capability matrix with the abduction column
% highlighted, above a compact schematic of the three operations on a trajectory.
\begin{figure}[t]
  \centering
  \begin{tikzpicture}[font=\small]
    % --- the capability matrix ---
    \matrix (m) [matrix of nodes,
        nodes={minimum width=2.5cm, minimum height=0.85cm, anchor=center,
               draw=gray!40, inner sep=3pt},
        column 1/.style={nodes={draw=none, text width=3.5cm, align=right}},
        row 1/.style={nodes={draw=none, font=\bfseries}},
        column sep=-\pgflinewidth, row sep=-\pgflinewidth]
    {
                                  & forward & control & abduction \\
      learned simulation (ours)   & \pass   & \pass   & \pass     \\
      frontier LLM (e.g.\ o3)     & \pass   & \pass   & \fail     \\
    };
    % shade + outline the abduction column (the gap, and the paper's subject)
    \begin{scope}[on background layer]
      \fill[orange!12] (m-1-4.north west) rectangle (m-3-4.south east);
    \end{scope}
    \draw[orange!75!red, very thick, rounded corners=1.5pt]
      (m-1-4.north west) rectangle (m-3-4.south east);

    % --- compact schematic of the three operations on a trajectory ---
    \begin{scope}[yshift=-2.55cm, font=\footnotesize]
      \node (x0) at (-3.0,0) {$x_0$};
      \node (xt) at (-0.2,0) {$x_T$};
      \draw[-{Latex[length=2mm]}, thick] (x0) -- (xt)
        node[midway, above=1pt] {forward: apply the rule};
      \node (g)  at (1.6,0) {goal};
      \draw[-{Latex[length=2mm]}, thick] (0.4,0) -- (g)
        node[midway, above=1pt] {control};
      % the inverse: a curved arrow running backward
      \node (obs) at (4.9,0) {noisy $x_{0:T}$};
      \node (xh)  at (3.0,-0.05) {$\hat{x}_0$};
      \draw[-{Latex[length=2mm]}, thick, orange!75!red]
        (obs) to[out=200,in=-20] (xh.east);
      \node[orange!75!red] at (4.0,-0.78) {\textbf{abduction}: recover the start};
    \end{scope}
  \end{tikzpicture}
  \caption{\textbf{The dissociation.} A reasoner can run a dynamical
  world \emph{forward}, \emph{control} it toward a target, or \emph{abduct} the initial
  condition from a noisy trajectory (the inverse). A frontier reasoning model (o3) closes
  the two forward-facing axes it shares with our learned-simulation reasoner, but not the
  inverse one. Our reasoner closes all three, sitting at the true-dynamics recoverability
  oracle. The orange column is the abduction gap.}
  \label{fig:dissociation}
\end{figure}
